%% ma: L12-B05-0022
%% lop: 12
%% chuong: 1
%% bai: 5
%% dang: 12.5.14
%% loai: TLN
%% muc: VD
%% dap_an: "3,82"
%% nguon: "(NBV)-ĐỀ SỐ 12-MỖI NGÀY 1 ĐỀ THI 2025.pdf (D12, MỖI NGÀY 1 ĐỀ - Đề số 12), Phần III Câu 3 (Sở Vĩnh Phúc 2025)"
%% kien_thuc: "Bài 5 (tối ưu bằng đạo hàm, hàm bậc ba), Bài 2 (giá trị lớn nhất, nhỏ nhất)"
%% trang_thai: chinh-thuc
%% ghi_chu: "Đã duyệt 13/10/2026. Đáp án gốc 3,82 đúng (sympy: L(x) = 16(44-3x)(6-x)(x-2), L'(x) = 0 tại x = (68 - 2 căn 283)/9 = 3,8172). Điều kiện không vượt quá 10 triệu không ràng buộc vì số ngày lưu trú dương kéo theo x < 6. Đã bỏ ảnh minh họa khu nghỉ dưỡng (không chứa dữ kiện toán)."
\begin{ex}
Một công ty kinh doanh dịch vụ nghỉ dưỡng nhận thấy rằng: Nếu áp dụng mức giá $3$ triệu đồng/người/ngày thì mỗi tháng có $140$ khách đến nghỉ và mỗi khách sẽ nghỉ $12$ ngày. Nếu cứ tăng giá thêm $500$ nghìn đồng/người/ngày thì hàng tháng số khách đến nghỉ sẽ giảm đi $6$ người và thời gian lưu trú của mỗi người khách cũng giảm $2$ ngày. Ngược lại, nếu cứ giảm $500$ nghìn đồng/người/ngày thì hàng tháng số khách đến nghỉ sẽ tăng thêm $6$ người và thời gian lưu trú của mỗi người khách cũng tăng thêm $2$ ngày. Hỏi công ty cần áp dụng mức giá bao nhiêu triệu đồng/người/ngày để lợi nhuận hàng tháng thu được là lớn nhất, biết tổng chi phí công ty phải chi cho một ngày lưu trú của mỗi người khách là $2$ triệu đồng và Sở du lịch không cho công ty thu vượt quá $10$ triệu đồng/người/ngày. (Làm tròn kết quả đến hàng phần trăm.)
\shortans{3,82}
\loigiai{Gọi $x$ (triệu đồng/người/ngày) là mức giá. Mỗi lần tăng $0{,}5$ triệu thì khách giảm $6$ người và ngày nghỉ giảm $2$ ngày, nên số khách là $140-12(x-3)=176-12x$ và số ngày nghỉ là $12-4(x-3)=24-4x$ (cần $2<x<6$ để lợi nhuận dương).

Lợi nhuận hàng tháng (triệu đồng):
\[L(x)=(176-12x)(24-4x)(x-2)=16(44-3x)(6-x)(x-2)=16\left(3x^3-68x^2+388x-528\right).\]
$L'(x)=16\left(9x^2-136x+388\right)=0\Leftrightarrow x=\dfrac{68\pm2\sqrt{283}}{9}$. Trong khoảng $(2;6)$ chỉ có $x_0=\dfrac{68-2\sqrt{283}}9\approx3{,}8172$; $L'>0$ trước $x_0$ và $L'<0$ sau $x_0$ nên $L$ lớn nhất tại $x_0$ (thỏa $x_0\le10$).

Vậy mức giá cần áp dụng xấp xỉ $3{,}82$ triệu đồng/người/ngày.}
\end{ex}
